% Source observation: analyses/018-2026-09-08-results-materials/observations/O007-kernel-realization.md
% Replacement prose for kernel-autoresearch.tex; Analysis 018 O007 and figure_data.json:runtime.
% Source reduction: evidence.py:runtime_subset; figure: plots.py:kernels.
% Underlying implementation audit: Run 029 observations/06-implementation-and-claim-audit.md.
\subsection{Realizing sparsity with specialized inference kernels}
\label{sec:kernel-autoresearch}

Logical sparsity becomes an execution benefit only if a kernel can skip work
at a lower cost than performing it. We evaluate 42 proposals from a
human-guided Codex search against native SDPA CUDA graphs on one RTX5090:
BF16, batch-one, uncached 2,048-token inference with full vocabulary logits.

The best qualified incumbent reaches $1.7830\times$ on the fixed A7-OL1,
$\kappa=0.5$ search checkpoint (Figure~\ref{fig:kernel-autoresearch}, left).
The selected kernel K050 qualifies on all 30 checkpoints in the
14M cohort, with geometric mean speedup $1.2340\times$. Speedup is positively
associated with sparsity (right; OLS $R^2=0.8167$), but checkpoint quality,
topology and weights vary; this is not an equal-quality acceleration comparison.

Relative to the same fusion with sparse paths disabled, K050 adds $31.1\%$
acceleration at the search checkpoint and $4.32\%$ geometrically across the
cohort, helping 14 of 30 checkpoints. Attention skipping is unprofitable:
disabling it improves every checkpoint, raising the mean to $1.2506\times$.
Appendix~\ref{app:kernel-results} details the implementation and ablations.
Whether longer sequences amortize attention-skip overhead remains untested.

Specialized kernels thus realize part of the logical opportunity, with
fusion and skip overhead determining the gain. Transfer to cached decoding,
larger models or other hardware is not established.

\begin{figure}[!htbp]
\centering
\includegraphics[width=\linewidth]{figures/07-kernel-realization.pdf}
\caption{\textbf{Sparse kernel auto-research for full-model acceleration.}
(a) Best qualified incumbent on fixed A7-OL1, $\kappa=0.5$, across 42
iterations, starting at native $1\times$ and reaching $1.7830\times$.
(b) K050 on 30 Pythia-14M checkpoints; symbols identify recipes. The dashed
line is a descriptive OLS fit with intercept ($R^2=0.8167$); dotted lines
mark native $1\times$. Speedup uses each checkpoint's native SDPA CUDA-graph
baseline on one RTX5090: BF16, batch one, uncached 2,048-token inference and
full vocabulary logits. Qualification covers all 338 validation blocks. All checkpoints
qualify, with geometric mean $1.2340\times$. Sparsity uses the canonical
FP16 measurement; both speedup axes have restricted ranges. The association
is not a causal sparsity effect. Appendix~\ref{app:kernel-results} gives
qualification details and ablations.}
\label{fig:kernel-autoresearch}
\end{figure}
\FloatBarrier
